\section{Trabalhos Relacionados}
\label{sec:relacionados}

\textbf{Posições a partir da fala.} A ciência política computacional estima posições de atores
políticos a partir das palavras que eles usam~\citep{laver2003extracting}, inclusive a de cada
parlamentar a partir do conjunto dos seus discursos~\citep{lauderdale2016measuring}.
\citet{izumi2021government} aplicam essa linha a discursos de senadores brasileiros, agregados por
partido. LLMs já foram usados para posicionar políticos tanto a partir do que o modelo sabe sobre o
nome~\citep{wu2023large} quanto a partir da leitura dos textos dos próprios
atores~\citep{lemens2025positioning}. Este trabalho substitui a escala por um perfil textual em
blocos, escrito somente com a fala de cada pessoa nas audiências de treino, e avalia o perfil pelo seu
uso na simulação. O efeito do nome observado por \citet{wu2023large} entra na condição 0, de
controle, que recebe só o nome e o cargo.

\textbf{Simulação de pessoas com LLMs.} \citet{argyle2023outofone} condicionam um LLM a perfis
sociodemográficos para reproduzir as respostas de subpopulações, e o condicionamento ao material do
próprio indivíduo supera o demográfico, seja com opiniões passadas no
prompt~\citep{hwang2023aligning}, seja com entrevistas longas~\citep{park2026llm}. Respostas de
personas, porém, variam com a redação do prompt e têm menos variância que as
humanas~\citep{bisbee2024synthetic}, e LLMs tendem a achatar e a estereotipar
identidades~\citep{wang2025large}. A simulação proposta trata de pessoas identificadas, usa como
material, além do nome e do cargo, apenas as falas públicas delas, declara a base da estimativa e é avaliada com opiniões publicadas sobre
audiências posteriores às do perfil.

\textbf{Atribuição.} A geração aumentada por recuperação~\citep{lewis2020rag} fundamenta o uso de
trechos das falas da pessoa, e a atribuição de um texto gerado a fontes
identificadas~\citep{rashkin2023measuring} costuma ser avaliada por anotadores ou por modelos, como
nas respostas com citações~\citep{gao2023enabling}. A UDV verifica por regra a forma do vínculo, sem
juiz automático, e a evidência é uma sentença localizada por offsets na fala da própria pessoa. Essa
verificação não avalia a escolha do codificador nem a sustentação de sentido.
